\chapter{Fundamentos teóricos y trabajos relacionados}
\label{chap:fundamentos}

\lettrine{E}{n} este capítulo se repasan los fundamentos técnicos más relevantes para la
comprensión del sistema desarrollado. Se presentan primero las técnicas de visión por
computador empleadas para detectar y seguir los vehículos, a continuación los conceptos de
sistemas distribuidos que sustentan la nueva arquitectura, y por último el mecanismo de
control de la velocidad mediante señales \gls{pwm}. El capítulo se cierra con una revisión
de los trabajos relacionados, en particular los dos Trabajos de Fin de Grado que constituyen
el punto de partida directo de este proyecto.

\section{Visión por computador}
\label{sec:vision}

La visión por computador, también denominada visión artificial, es un área multidisciplinar que tiene como propósito dotar a los sistemas computacionales de la capacidad de percibir, procesar y comprender su entorno a partir de imágenes o secuencias de vídeo. Según G. Bradski y A. Kaehler~\cite{bradski2008}, esta disciplina se fundamenta en la transformación de datos visuales crudos en decisiones ejecutables o en nuevas representaciones geométricas y semánticas del entorno. 

Aunque para los seres vivos la procesar estímulos visuales es casi instantáneo. Replicar estas capacidades en sistemas artificiales constituye un reto altamente complejo. Su desarrollo combina principios de matemáticas, física, procesamiento de señales, ciencias de la computación, ... Impulsado en gran medida por librerías de propósito general como OpenCV~\cite{opencv}, que han estandarizado y acelerado el procesamiento en tiempo real.

En el alcance de este proyecto, la visión por computador se usa como el sensor principal del sistema. Permite capturar imagenes para luego procesarlas con el objetivo de identicar los coches para su posterior control usando esta inforamción. Para conseguirlo se emplean un conjunto de técnicas:

\begin{itemize}
    \item \textbf{Segmentación por color en espacio HSV:} Este tipo de espacio de color es mucho más robusto a los cambios  de iluminación, sobras y reflejos. Lo que facilita detectar los colores específicos asociados a pegatinas.     
    \item \textbf{Extracción de contornos y calculo de centros geometricos} 
    \item \textbf{Seguimiento de 2 marcadores:} Sistema de posicionamiento basado en 2 pegatinas, una frontal y una trasera, permiten controlar la posición del coche y el angulo del derrape.    
    \item \textbf{Regiones de interés (ROI) y máscaras de trayectoria:} Delimitación del area de búsqueda sobre el trazado del circuito para reducir la carga computacional y eliminar falsas detecciones en los márgenes de la imagen.
\end{itemize}

\subsection{Seguimiento de objetos}
\label{sec:seguimiento}

El seguimiento de objetos, o \textit{tracking}, es el proceso algorítmico que permite obtener la trayectoria de los objetos en una secuencia de imagenes~\cite{bradski2008}. En lugar de tratar cada imagen de manera independiente, los algoritmos de seguimiento combinan la detección visual instantánea con modelos de predicción basados en el estado previo del objeto.

Para garantizar robustez frente a perturbaciones, es crítico seleccionar las propiedades visuales más distintivas y estables del objetivo. Normalmente las características más usadas son:

\begin{itemize}
    \item \textbf{Color:} permite discriminar de forma muy sencilla y eficiente para aislar un objeto del fondo. El problema es lo dependiente que somos a los cambios de la luz en le entorno, ya que nos podemos ver muy afectados.
    \item \textbf{Bordes y contornos:} Representan las zonas con saltos de alta frecuencia, estos cambios se pueden identificar lo que permite identifcar la forma del objeto.    
    \item \textbf{Flujo óptico:} Se usa un campo vectorial que describe el movimiento aparente de los píxeles entre fotogramas consecutivos. Permite identificar el patrón de desplazamiento del objeto. 
    \item \textbf{Textura:} Cuantifica los patrones de rugosidad o contraste en la superficie de un objeto.
\end{itemize}

En escenarios reales, todo sistema de \textit{tracking} debe enfrentarse a factores que degradan la estimación, tales como oclusiones parciales o totales, desenfoque por movimiento a altas velocidades (\textit{motion blur}) o variaciones abruptas de iluminación artificial.

\subsection{Región de interés (ROI) y extrapolación}
\label{sec:roi}

La región de interés (\gls{roi}) es un subconjunto de la imagen, generalmente rectangular,
que delimita el área donde se espera encontrar la información relevante. Procesar únicamente
la \gls{roi} en lugar del fotograma completo permite reducir el coste computacional y disminuir el ruido al excluir regiones irrelevantes donde podrían producirse falsas detecciones.

Para situar la \gls{roi} en cada fotograma se recurre a la extrapolación lineal: conocidas
las dos últimas posiciones del coche, $(x_{n-1}, y_{n-1})$ y $(x_n, y_n)$, la posición
siguiente se estima:
\[
    x_{n+1} \approx x_n + (x_n - x_{n-1}), \qquad
    y_{n+1} \approx y_n + (y_n - y_{n-1}),
\]
y la ventana de búsqueda se centra en el punto predicho. Cuando la detección falla (por
ejemplo, por una oclusión o porque el coche abandona el campo de visión de la cámara), la
ventana se amplía progresivamente en los fotogramas siguientes hasta cubrir la imagen
completa, garantizando la readquisición del vehículo en cuanto vuelve a ser visible.

\section{PWM}
La modulación por ancho de pulso (\gls{pwm}) es una técnica basada en la generación de una señal digital cuadrada, constituida por un tren de pulsos de amplitud fija y frecuencia constante, cuya longitud es variable. Esta variación viene definida por el ciclo de trabajo o \textit{duty cycle}, que representa la proporción de tiempo que la señal permanece en nivel alto respecto a su período total. Aunque la señal solo alterna entre dos valores discretos (0\,V y la tensión de alimentación), el sistema o la carga conectada percibe una tensión media analógica directamente proporcional a dicho ciclo de trabajo. De este modo, se logra emular un control continuo de potencia. Un caso de uso muy común para controlar la velocidad de los motores. 

\section{Sistemas distribuidos}
\label{sec:sistemas-distribuidos}

Un sistema distribuido es un conjunto de computadores independientes que se presenta ante
sus usuarios como un único sistema coherente. Sus principales ventajas son
la escalabilidad --la capacidad de crecer añadiendo más máquinas-- y la tolerancia a fallos
--la posibilidad de que el sistema continúe funcionando aunque alguno de sus componentes
falle--. A cambio, introduce retos que no existen en las aplicaciones monolíticas: los
procesos no comparten memoria y deben coordinarse mediante mensajes, la red añade latencia y
puede perder paquetes, y los fallos son parciales, es decir, una parte del sistema puede
caer mientras el resto sigue en marcha.

Este trabajo convierte el sistema monolítico heredado en un sistema distribuido
precisamente para ganar escalabilidad: el número de cámaras deja de estar limitado por los
puertos y la capacidad de procesamiento de un único ordenador. Los conceptos de esta sección
fundamentan las decisiones de diseño tomadas para afrontar los retos mencionados.

\subsection{Paradigma publicador/suscriptor}
\label{sec:pubsub}

En el modelo publicador/suscriptor, los procesos no se comunican directamente entre sí, sino
a través de canales con nombre (\textit{topics}): los publicadores envían mensajes a un
topic sin conocer quién los recibirá, y los suscriptores expresan su interés en un topic sin
conocer quién publica en él. Este desacoplamiento aporta grandes ventajas que simplican el diseño. 

El paradigma encaja de forma natural en el problema de este proyecto: cada nodo cámara
publica las posiciones que detecta sin saber cuántos controladores existen, y cada
controlador se suscribe a las posiciones de su coche sin saber cuántas cámaras cubren el circuito. Añadir una cámara o un coche no requiere reconfigurar ningún otro componente del
sistema.

\subsection{Middleware DDS y calidad de servicio (QoS)}
\label{sec:dds-qos}

\gls{dds} es un estándar de middleware publicador/suscriptor orientado a sistemas de tiempo
real, definido por el Object Management Group~\cite{ddsspec}, sobre el que ROS2 construye
toda su capa de comunicación. Dos de sus características son especialmente
relevantes para este trabajo. La primera es el descubrimiento automático: los participantes
de un mismo dominio (identificado por un número, el \textit{domain ID}) se anuncian
periódicamente en la red local y establecen las conexiones entre publicadores y suscriptores
sin ninguna configuración manual de direcciones. La segunda es la calidad de servicio
(\gls{qos}): cada publicador y suscriptor declara sus garantías de entrega mediante un
conjunto de políticas, entre las que destacan:

\begin{itemize}
    \item \textbf{Fiabilidad (\textit{reliability}):} en modo \textit{reliable} el
    middleware reenvía los mensajes perdidos hasta confirmar su entrega; en modo
    \textit{best effort} los mensajes se envían una única vez y pueden perderse.

    \item \textbf{Durabilidad (\textit{durability}):} en modo \textit{volatile} un
    suscriptor solo recibe los mensajes publicados después de conectarse; en modo
    \textit{transient local} el publicador conserva los últimos mensajes y los entrega a
    los suscriptores que se conecten más tarde.

    \item \textbf{Historial (\textit{history} y \textit{depth}):} determina cuántos
    mensajes se almacenan en las colas de envío y recepción.
\end{itemize}

La elección de estas políticas supone un compromiso entre latencia y garantías de entrega.
Para los datos de posición, que se publican decenas de veces por segundo, se emplea el
perfil de datos de sensor (\textit{best effort}): un mensaje perdido carece de importancia
porque enseguida llega el siguiente, y resulta preferible perder un fotograma a acumular
retardo reenviando información ya obsoleta. En cambio, para la posición de la línea de meta,
que se publica una única vez al arrancar, se utiliza un perfil \textit{reliable} y
\textit{transient local}: el mensaje debe llegar con garantías incluso a los controladores
que arranquen después de que la cámara lo haya publicado.

%\subsection{Tolerancia a fallos: heartbeat, failover y watchdog}
%\label{sec:tolerancia-fallos}
%
%En un sistema distribuido, la caída de un proceso no puede detectarse de forma directa: la
%única evidencia disponible es la ausencia de comunicación. Sobre esta idea se construyen los
%tres mecanismos de tolerancia a fallos empleados en este proyecto:
%
%\begin{itemize}
%    \item \textbf{Latido de vida (\textit{heartbeat}):} un proceso emite periódicamente un
%    mensaje que atestigua que sigue vivo. Quien lo vigila considera que ha fallado cuando
%    transcurre un plazo sin recibir latidos.
%
%    \item \textbf{Replicación primario/respaldo (\textit{failover}):} un componente crítico
%    se despliega por duplicado. La réplica primaria realiza el trabajo y emite latidos; la
%    réplica de respaldo procesa la misma información en silencio y, si detecta la ausencia
%    de latidos, asume el rol de primaria de forma automática. De este modo el servicio se
%    recupera en segundos sin intervención humana.
%
%    \item \textbf{Perro guardián (\textit{watchdog}):} un componente vigila que otro siga
%    enviando información con la cadencia esperada y, en caso contrario, ejecuta una acción
%    de seguridad. En este proyecto, el nodo puente detiene la alimentación de un carril
%    cuando deja de recibir consignas de velocidad para él, de forma que un coche cuyo
%    control ha fallado (o que se ha salido de la pista) queda detenido sin afectar al resto
%    de la carrera.
%\end{itemize}
%
%Como tercera capa de recuperación, el sistema de lanzamiento de ROS2 permite reiniciar
%automáticamente (\textit{respawn}) cualquier nodo que termine de forma inesperada, de manera
%que un nodo caído se reincorpora al sistema al cabo de unos segundos.

\section{Control por Modulación por Ancho de Pulso (PWM)}
\label{sec:pwm-fundamentos}

La modulación por ancho de pulso (\gls{pwm}) es una técnica que genera una señal cuadrada de
amplitud fija que alterna entre 0\,V y la tensión de alimentación a una frecuencia
constante, variando únicamente la proporción de tiempo en la que permanece en nivel alto,
denominada ciclo de trabajo o \textit{duty cycle}. Aunque la salida es
digital, la carga percibe una tensión media equivalente a la tensión de alimentación
multiplicada por el ciclo de trabajo, emulando así un control analógico de potencia con muy
pocas pérdidas, ya que los transistores de conmutación operan siempre completamente
encendidos o apagados.

\section{Trabajos relacionados}
\label{sec:traballos-relacionados}

El control automático de vehículos en pistas de Scalextric es un problema con una larga
trayectoria en entornos didácticos y de prototipado, abordado tanto por soluciones
comerciales como por proyectos académicos. Esta sección se centra en la
línea de TFGs en los que se cimienta este trabajo.

\subsection{Línea de trabajos previos}

La línea comienza con el TFG de Miguel Fernández López, centrado en la detección y el
seguimiento visual de un coche de Scalextric, incluyendo la identificación de sectores de la
pista y el cálculo de tiempos de vuelta, y con el TFG de Julián Barcia Facal, que aportó el
diseño del hardware de control basado en Arduino y señales \gls{pwm}. Ambos trabajos están
descritos en detalle en las memorias posteriores~\cite{tfgmario,tfgadrian}.

\subsection{TFG de Mario López Cea}
\label{sec:tfg-mario}

El trabajo de Mario López Cea~\cite{tfgmario} logró por primera vez en esta línea el control
autónomo y adaptativo del coche. Con una única cámara fija sobre el circuito, el sistema
detecta dos etiquetas de color colocadas en el coche (delantera y trasera), reconstruye su
trayectoria, calcula su velocidad instantánea y detecta los derrapes. Sobre esta base diseñó
dos algoritmos de control de velocidad: uno de maximización de la velocidad media y otro de
anticipación a las zonas de derrape, que aprende las curvas peligrosas del circuito y frena
el coche antes de alcanzarlas.

De este trabajo se heredan directamente tres elementos. El primero es el esquema de doble
pegatina, que este proyecto mantiene como base de la detección, así como las tecnicas de evaluación de
tiempos por etapas del procesamiento, que en este proyecto se replica de forma distribuida
mediante mensajes de telemetría.

\subsection{TFG de Adrián Rego Criado}
\label{sec:tfg-adrian}

El trabajo de Adrián Rego~\cite{tfgadrian} extendió el sistema a circuitos grandes
sustituyendo la cámara única por una red de cámaras con campos de visión solapados o no, que
cubren la totalidad del trazado. Introdujo además la detección automática de los coches y de
la línea de meta mediante segmentación por color (eliminando la configuración manual), y una interfaz gráfica organizada en paneles.

Su limitación fundamental, señalada explícitamente en su propia memoria como línea de
trabajo futuro, es que toda la aplicación se ejecuta como un único proceso en un único
ordenador, al que deben conectarse físicamente todas las cámaras. De este trabajo se heredan la detección automática de la línea de meta por color y
la filosofía de calibración automática sin configuración manual, todas ellas adaptadas a la
arquitectura de nodos distribuidos.

\subsection{Posicionamiento de este trabajo}
\label{sec:posicionamento}

Este proyecto recoge el testigo de ambos trabajos y aborda precisamente la evolución que el
segundo dejaba planteada: convertir el sistema en una arquitectura distribuida. Los
principales avances respecto a los trabajos de partida son:

\begin{itemize}
    \item Descomposición de la aplicación monolítica en nodos ROS2 independientes,
    ejecutables en máquinas distintas de la red local.

    \item Escalabilidad horizontal: el número de cámaras deja de estar limitado por un
    único equipo; añadir cobertura al circuito se reduce a lanzar un nuevo nodo cámara.

    \item Adaptación del algoritmo de control de velocidad al entorno multicámara.

    \item Tolerancia a fallos con reinicio automatico de los nodos.

    \item Despliegue reproducible mediante contenedores Docker, sin necesidad de instalar
    ROS2 en las máquinas.

    \item Telemetría distribuida y herramienta de grabación y análisis de las sesiones.
\end{itemize}
